% Generated draft: requires mathematical and source review.
\documentclass[11pt]{article}
\usepackage[margin=25mm]{geometry}
\usepackage{fontspec}
\setmainfont{Latin Modern Roman}
\usepackage{amsmath,amssymb,mathtools,unicode-math}
\setmathfont{Latin Modern Math}
\usepackage{enumitem,needspace,xurl,hyperref}
\setlength{\parindent}{0pt}
\setlength{\parskip}{5pt}
\setlength{\emergencystretch}{3em}
\newcommand{\N}{\mathbb{N}}
\newcommand{\Z}{\mathbb{Z}}
\newcommand{\Q}{\mathbb{Q}}
\newcommand{\R}{\mathbb{R}}
\newcommand{\C}{\mathbb{C}}
\newcommand{\F}{\mathbb{F}}
\newcommand{\A}{\mathbb{A}}
\newcommand{\PP}{\mathbb{P}}
\newcommand{\E}{\mathbb{E}}
\newcommand{\eps}{\varepsilon}
\newcommand{\dd}{\,\mathrm{d}}
\begin{document}
\section*{UnsolvedMath ID: 700594}
\textbf{Source problem code:} catalogue entry
\subsection{Definitions and mathematical statement}
A \emph{variety of groups} is the class of all groups satisfying some specified set of group identities. Let \(\mathfrak V\) denote the set of all varieties of groups. Write \(M(\mathcal V)\) provisionally for the Mal'cev theory of \(\mathcal V\) in the sense of the cited work of D. M. Smirnov. The supplied record defines
\[
\mathcal V\sim_S\mathcal W
\quad\Longleftrightarrow\quad
M(\mathcal V)=M(\mathcal W).
\]
For \(\mathcal V\in\mathfrak V\), its full \(S\)-equivalence class would therefore be
\[
[\mathcal V]_S
=\{\mathcal W\in\mathfrak V:M(\mathcal W)=M(\mathcal V)\}.
\]
\textbf{Source question (D. M. Smirnov, 1984).} What is the cardinality of the set of \(S\)-equivalent varieties of groups?

\textbf{Unresolved qualifications.} The supplied record does not define Mal'cev theory beyond referring to Smirnov's article, so \(M\) above is a placeholder for the source-specific invariant, not an independently specified theory. If the intended set is a full equivalence class, the requested quantity is \(\lvert[\mathcal V]_S\rvert\); however, the record neither specifies \(\mathcal V\) nor states whether the question asks for a common cardinality, the possible cardinalities, or a maximal cardinality. It also does not explicitly require the family of pairwise \(S\)-equivalent varieties to be a full equivalence class. These choices remain unresolved.
\subsection{Short English statement}
Smirnov asks how many varieties of groups can have the same Mal'cev theory. A precise cardinality question requires the cited definition of that theory and clarification of which equivalence class or family is intended. The supplied registry labels the problem open; this is not an independently verified status assessment.
\subsection{Sources}
\begin{itemize}
\item Local UnsolvedMath registry, supplied complete source record: permanent ID \(700594\), exact source problem code \texttt{catalogue entry}, title \emph{Kourovka Notebook Problem 9.61}; upstream revision \texttt{d371ee917ac9647fbb06cd6fc529deec6f3dda82}. The record supplies the status label \emph{open}.
\item \emph{The Kourovka Notebook (Unsolved Problems in Group Theory)}, problem 9.61, proposed by D. M. Smirnov in the 9th issue (1984). The supplied background identifies the 21st issue edition (2026), \texttt{arXiv:1401.0300v48}; the collection description names E. I. Khukhro and V. D. Mazurov as editors.
\item D. M. Smirnov, \emph{Algebra and Logic}, volume 22, no.\ 6 (1983), pp.\ 492--501, cited in the supplied statement for Mal'cev theory. No article title or further definition is supplied. This citation and the notebook attribution have not been independently checked.
\end{itemize}
\end{document}
